\chapter{Sistema de predicción: agentes y meta-agente}
\label{chap:agentes}

La hipótesis de partida no es que una sola familia de ratios domine de forma estable todos los
regímenes, sino que distintas vistas económicas del universo aportan información con persistencia
desigual. De ahí los cinco agentes especializados ---calidad, valor, crecimiento, tendencia y
riesgo--- y un combinador que aprende a ponderarlos sin saber de antemano cuál será el más útil.
Este capítulo describe ese mecanismo y las alternativas que el catálogo ofrece en cada punto; qué
valores se eligieron finalmente depende de una regla de decisión que se explica en el capítulo
siguiente.

\section{Arquitectura y aprendizaje \textit{walk-forward}}

Cada agente recibe sólo los bloques de variables asignados a su especialidad. El de calidad resume
rentabilidad, márgenes, eficiencia y solidez financiera; el de valor utiliza ratios de precio y
flujos; el de crecimiento mide aceleración y estabilidad de fundamentales; el de tendencia recoge
momentum de precio y de fundamentales; y el de riesgo representa riesgo de precio y liquidez. Esta separación tiene dos fines: conserva interpretabilidad económica y evita que el
meta-agente confunda cinco copias del mismo predictor.

En cada fecha de snapshot se ajusta un LightGBM por agente sobre una ventana móvil de ocho años. El
catálogo declara la rejilla de hiperparámetros --- profundidad, número de estimadores, tasa de
aprendizaje y mínimo de observaciones por hoja --- y el protocolo del capítulo siguiente decide qué
valor toma cada uno. LightGBM permite representar interacciones no lineales sin imponer una forma
funcional única (Ke et al., 2017); el entrenamiento temporal evita que el modelo evalúe
observaciones que habría visto durante el ajuste.

Las cinco puntuaciones se convierten en rangos transversales y pasan al meta-agente, que estima una
Ridge positiva sobre cohortes trimestrales ya cerradas. Mientras no existan etiquetas cerradas
suficientes, el meta reparte por construcción pesos iguales de 0,20.

El catálogo ofrece dos variantes de esa combinación, y elegir entre ellas es una hipótesis de diseño,
no un detalle de implementación. La variante \textbf{acotada} obliga a que ningún agente pese menos
del 10\,\% ni más del 50\,\%: representa la idea de que una arquitectura de cinco especialistas sólo
tiene sentido si ninguno puede apagar a los demás. La variante \textbf{libre} no impone tope alguno,
y representa la idea contraria: si un especialista domina de forma consistente, lo correcto es dejar
que el procedimiento converja hacia él.

Las dos son defendibles y miden cosas distintas. La acotada protege una propiedad cualitativa del
diseño ---que sigan siendo cinco agentes--- que ninguna métrica del protocolo evalúa; la libre
maximiza el criterio que el protocolo sí mide. Por eso la variante figura en el catálogo cerrado
como una variable más, y no como una preferencia del autor.

Conviene fijarse en que el meta no ve las puntuaciones sino los \emph{rangos}: es lo que permite
combinar agentes cuyas escalas no son comparables entre sí.

\subsection{Qué predice el modelo}

El sistema no aprende a predecir cuánto subirá una acción, sino en qué puesto quedará. Para cada
fecha se toma el retorno a doce meses de todas las acciones cuya etiqueta ya puede cerrarse, se
ordenan de peor a mejor y se convierte esa posición en un número entre 0 y 1:
\[
 y_{i,t}=\frac{\operatorname{rank}_{j\in\mathcal U_t}
 \left(R_{j,t\rightarrow t+12}\right)-1}{|\mathcal U_t|-1}.
\]
Un valor de 1 identifica la mejor acción de su fecha y un 0 la peor. Aprender el puesto en vez del
nivel tiene dos ventajas: alinea el entrenamiento con la métrica que después decide el ganador, y
hace al sistema inmune a que un año entero suba o baje. No elimina, en cambio, la dependencia entre
acciones que comparten fecha, que es la razón de agregar toda la inferencia por cohorte.

Cada agente $a$ estima una función no lineal
$\hat y^{(a)}_{i,t}=f_a(X^{(a)}_{i,t};\theta_{a,t})$, donde
$X^{(a)}_{i,t}$ contiene únicamente las variables permitidas a ese especialista y disponibles en
$t$. El subíndice temporal de $\theta$ importa: no existe un único árbol ajustado una vez sobre toda
la historia, sino una sucesión de modelos. Una predicción publicada en $t$ queda ligada a la versión
entrenada con la última información cerrada antes de $t$.

\subsection{Dos relojes que nunca se confunden}

La implementación mantiene separados el reloj con el que se emite una señal y el reloj con el que se
aprende de ella. En una fecha dada puede puntuarse una empresa con información ya observable, pero
esa misma fila no puede entrar en un entrenamiento hasta que transcurran los doce meses que hacen
falta para conocer su resultado. Eso impide una fuga menos visible que usar un precio futuro: que
una fila reciente, cuyas características se conocen pero cuyo retorno no, intervenga en la selección
de variables o en el meta. El cierre de cohortes es por tanto una condición previa del
entrenamiento, la poda y la calibración, no una corrección posterior.

\section{Entrenamiento y combinación de rangos}

La especialización no asigna una conclusión previa al modelo, sino un vocabulario económico con el
que formular hipótesis: el agente de calidad recibe ratios como ROE, ROIC, márgenes o endeudamiento,
pero no la regla «ROIC alto implica comprar», que tendría que aprender de la historia disponible. La
amplitud de ese vocabulario es también una variable del catálogo, que ofrece un conjunto reducido de
bloques o los once disponibles; en ambos casos una poda por estabilidad fuera de muestra limita
cuántas variables conserva cada agente. Un conjunto demasiado estrecho amputa información y uno
demasiado amplio añade varianza, así que la elección se mide en vez de decidirse a priori.

En cada snapshot el modelo recibe sólo la historia permitida. Si la fecha es abril de 2020 y la
ventana es de ocho años, el conjunto comienza aproximadamente en abril de 2012; una fila de finales
de 2019 entra sólo si su etiqueta anual ya cerró. El ajuste se repite al avanzar el tiempo. Los
hiperparámetros que gobiernan la complejidad del árbol --- profundidad, número de estimadores, tasa
de aprendizaje y mínimo de observaciones por hoja --- controlan cuán específicas pueden ser las
interacciones aprendidas, y sus valores no se fijan aquí por criterio general sino con la regla del
Capítulo~\ref{chap:diseno-experimental}.

Los scores se convierten a rangos porcentuales dentro de cada cohorte. Esta operación evita que la
escala numérica de un agente domine al resto: si el de riesgo produce valores entre 0,4 y 0,6 y el
de valor entre $-2$ y $3$, ambos pasan a expresar posición relativa en su sección transversal.
El meta combina después esos rangos, no sus unidades originales.

Antes del ajuste, los extremos se tratan con reglas declaradas y los valores ausentes se imputan
usando únicamente estadísticas obtenidas en el corte de entrenamiento. Una variable no se conserva
porque explique retrospectivamente toda la muestra: su disponibilidad y su contribución deben
sostenerse en los pliegues temporales permitidos. Con ello, la poda reduce dimensión sin convertir
el conjunto reservado en un selector informal. El diccionario completo de variables, direcciones,
bloques y agentes figura en el Apéndice~\ref{anexo:catalogo}.

La pregunta que el meta-agente resuelve en cada fecha es sencilla de enunciar: qué combinación de
los cinco rangos habría acertado mejor el orden real en las cohortes que ya cerraron. Formalmente,
busca los pesos $\boldsymbol w$ que minimizan el error de esa combinación, con un término adicional
que penaliza los pesos grandes:
\[
 \widehat{\boldsymbol w}_t=
 \arg\min_{\boldsymbol w}
 \sum_{(i,s)\in\mathcal C_t}
 \left(y_{i,s}-w_0-\sum_{a=1}^{5}w_a z^{(a)}_{i,s}\right)^2
 +\lambda\lVert\boldsymbol w\rVert_2^2,
\]
donde $z^{(a)}_{i,t}$ son los cinco rangos y $\mathcal C_t$ contiene sólo cohortes cerradas
anteriores a $t$. Los pesos no pueden ser negativos y se normalizan para sumar uno; la variante
acotada añade además que ninguno baje de 0,10 ni pase de 0,50.

Ese término de penalización tiene una función concreta: cuando dos agentes dicen cosas parecidas,
impide que sus pesos salten bruscamente de uno a otro ante cambios pequeños en los datos. Lo que no
puede evitar es que, si un agente resulta consistentemente mejor, la variante libre acabe
concentrándose en él. Eso es una hipótesis que el protocolo contrasta, no un defecto a corregir.

Cuando $\mathcal C_t$ todavía no tiene la longitud mínima exigida, no se ajusta una regresión
frágil. El sistema utiliza el reparto igualitario y registra el estado de arranque. Esta conducta es
parte del contrato reproducible: dos ejecuciones con el mismo catálogo y la misma instantánea no
pueden decidir de manera distinta si ya existe historia suficiente.

\section{Ejemplo de combinación}

\begin{ejemplo}[De cinco rangos a un score]
Una empresa presenta rangos quality 0,80, value 0,55, growth 0,70, momentum 0,45 y risk 0,90. En el
arranque, los pesos iguales producen $(0,80+0,55+0,70+0,45+0,90)/5=0,68$: queda algo por encima de
la mediana del universo, y ninguna de sus cinco vistas pesa más que otra.

Si la evidencia ya cerrada indicase que el agente de riesgo ordena mejor que los demás y el meta le
asignara 0,70, repartiendo 0,075 entre los cuatro restantes, el score pasaría a
$0{,}70\times0{,}90 + 0{,}075\times(0{,}80+0{,}55+0{,}70+0{,}45) = 0{,}82$.
\end{ejemplo}

Los rangos del ejemplo son inventados, pero el cálculo es el real, y muestra qué está en juego entre
las dos variantes. Con la acotada, ese peso de 0,70 sería imposible: el tope del 50\,\% limita
cuánto puede desplazarse la puntuación hacia el especialista dominante. Con la libre nada lo impide
mientras la evidencia lo respalde, hasta el extremo de que la puntuación final acabe siendo, en la
práctica, la de un solo agente.

Conviene delimitar qué permite afirmar esta arquitectura. Los pesos hacen \emph{atribuible} el
resultado a agentes concretos, pero no permiten inferir causalidad: que un especialista reciba peso
alto no demuestra que su lectura económica sea la causa del retorno, y puede reflejar interacciones
entre variables o límites de los datos. Por eso el trabajo no se detiene en los pesos y añade
después estabilidad por eras, neutralización por estilo y atribución local dentro de cada agente.
